\subsection{Using a Calculator to Simplify} 
Sometimes, when we simplify a root, we need to search many possible factors.
\begin{example}Simplify $\sqrt{108}$.
\begin{lpic}{}{8}
\foreach \y/\ys/\q/\e in {1/1/108/=,
													2/4/27/=,
													3/9/12/=,
													4/16/6.75/=,
													5/25/4.32/=,
													6/36/3/=,
													7/49/2.204/\approx,
													8/64/1.6875/=} {
	\regtextleft{0.25}{-\y}{$\y^2=\ys$};
	\regtextleft{4.25}{-\y}{$108\div\ys\e\q$};
}
\regtextleft{11.25}{-1}{$\sqrt{108}={}$};
\end{lpic}
\end{example}
We can use a computer to do the repetitive arithmetic. The procedure here uses the TI-84 calculator, but it can be done with many other calculators or with a computer.
\begin{procedure}{Simplifying Roots with the TI-84}
\begin{enumerate*}
\item Press \TIbut{y=} to access the function entry screen.
\item Enter \TI{X\TIsup{p}} as \TI{Y\TIsub{1}} and \TI{\#\#/X\TIsup{p}} as \TI{Y\TIsub{2}}, where \TI{\#\#} is the number whose root you want to simplify, and \TI{p} is the index.
\item Press \TIbut{2nd}\ \TIbut{graph} to access the ``Table'' feature.
\item The \TI{Y\TIsub{1}} column shows the perfect $p$th powers, the \TI{X} column shows the $p$th root, and the \TI{Y\TIsub{2}} column shows the remaining factor. You are looking for the largest possible \TI{X} which gives an integer in the \TI{Y\TIsub{2}} column. You know you're done searching when the number in the \TI{Y\TIsub{2}} is 2 or less.
\end{enumerate*}
\end{procedure}
\begin{example} Simplify $\sqrt[3]{112}$.
\begin{lpic}{}{3}
\end{lpic}
\end{example}